\subsection{BÀI TẬP TRẮC NGHIỆM}

\ind{PHẦN I.} \inden{Câu trắc nghiệm nhiều phương án lựa chọn. Mỗi câu hỏi học sinh chỉ chọn một phương án.}\\
\setcounter{ex}{0}
\Opensolutionfile{ans}[ans/1D3-B2-1]
\begin{ex}%[1D1Y2-2]
	Trong các số đặc trưng đo xu thế trung tâm dưới đây, số nào thỏa mãn có $25\% $ giá trị trong mẫu số liệu nhỏ hơn nó và $75\% $ giá trị trong mẫu số liệu lớn hơn nó?
	\choice
	{số trung bình}
	{trung vị}
	{\True tứ phân vị thứ nhất}
	{tứ phân vị thứ ba}
	\loigiai{
		Tứ phân vị thứ nhất $Q_1$ có $25\% $ giá trị trong mẫu số liệu nhỏ hơn $Q_1$ và $75\% $ giá trị trong mẫu số liệu lớn hơn $Q_1$.}
\end{ex}

\begin{ex}%[1D1Y2-2]
	Trong các số đặc trưng đo xu thế trung tâm dưới đây, số nào thỏa mãn có $75\% $ giá trị trong mẫu số liệu nhỏ hơn nó và $25\% $ giá trị trong mẫu số liệu lớn hơn nó?
	\choice
	{số trung bình}
	{trung vị}
	{tứ phân vị thứ nhất}
	{\True tứ phân vị thứ ba}
	\loigiai{
		Tứ phân vị thứ ba $Q_3$ có $75\% $ giá trị trong mẫu số liệu nhỏ hơn $Q_3$ và $25\% $ giá trị trong mẫu số liệu lớn hơn $Q_3$.}
\end{ex}

\begin{ex}%[1D1B2-1]
	Khảo sát thời gian tập thể dục của một số học sinh khối $11$ thu được mẫu số liệu ghép nhóm sau
	\begin{center}
		\begin{tabular}{|c|c|c|c|c|c|}
			\hline Thời gian (phút) & {$[0 ; 20)$} & {$[20 ; 40)$} & {$[40 ; 60)$} & {$[60 ; 80)$} & {$[80 ; 100)$} \\
			\hline Số học sinh & $5$ & $9$ & $12$ & $10$ & $6$ \\
			\hline
		\end{tabular}
	\end{center}
	Nhóm chứa trung vị của mẫu số liệu trên là
	\choice{ \True $[40 ; 60)$}
	{ $[20 ; 40)$}
	{ $[60 ; 80)$}
	{ $[80 ; 100)$}
	\loigiai{
		Ta có $n=42$.\\
		Nên trung vị của mẫu số liệu trên là $Q_2=\dfrac{x_{21}+x_{22}}{2}$.\\
		Mà $x_{21}$, $ x_{22} \in[40 ; 60)$.\\
		Vậy nhóm chứa trung vị của mẫu số liệu trên là nhóm $[40 ; 60)$.}
\end{ex}

\begin{ex}%[1D1B2-2]
	Khảo sát thời gian tập thể dục của một số học sinh khối $11$ thu được mẫu số liệu ghép nhóm sau
	\begin{center}
		\begin{tabular}{|c|c|c|c|c|c|}
			\hline Thời gian (phút) & {$[0 ; 20)$} & {$[20 ; 40)$} & {$[40 ; 60)$} & {$[60 ; 80)$} & {$[80 ; 100)$} \\
			\hline Số học sinh & $5$ & $9$ & $12$ & $10$ & $6$ \\
			\hline
		\end{tabular}
	\end{center}
	Nhóm chứa tứ phân vị thứ nhất của mẫu số liệu trên là
	\choice{ \True$[40 ; 60)$}
	{ $[20 ; 40)$}
	{ $[60 ; 80)$}
	{ $[80 ; 100)$}
	\loigiai{
		Ta có $n=42$.\\
		Nên tứ phân vị thứ nhất của mẫu số liệu trên là $Q_1=x_{11}$.\\
		Mà $x_{11} \in[20 ; 40)$.\\
		Vậy nhóm chứa tứ phân vị thứ nhất của mẫu số liệu trên là nhóm $[20 ; 40)$.
	}
\end{ex}

\begin{ex}%[1D1B2-2]
	Khảo sát thời gian tập thể dục của một số học sinh khối $11$ thu được mẫu số liệu ghép nhóm sau
	\begin{center}
		\begin{tabular}{|c|c|c|c|c|c|}
			\hline Thời gian (phút) & {$[0 ; 20)$} & {$[20 ; 40)$} & {$[40 ; 60)$} & {$[60 ; 80)$} & {$[80 ; 100)$} \\
			\hline Số học sinh & $5$ & $9$ & $12$ & $10$ & $6$ \\
			\hline
		\end{tabular}
	\end{center}
	Tứ phân vị thứ ba của mẫu số liệu trên là
	\choice{$115$}
	{$117$}
	{\True  $71$}
	{$112$}
	\loigiai{
		Số học sinh được khảo sát là $n=5+9+12+10+6=42$.\\
		Gọi $x_1 ; x_2 ; \ldots ; x_{42}$ là thời gian tập thể dục của $42$ học sinh lớp $11$ được xếm theo thứ tự không giảm.
		\begin{align*}
			& Q_2=\dfrac{1}{2}\left(x_{21}+x_{22}\right) \in[40 ; 60) \\
			& Q_1=x_{11} \in[20 ; 40) \\
			& Q_3=x_{31} \in[60 ; 80) \\
			& x_i=60, c=26, n=42, n_j=10\\
			&Q_3=60+\dfrac{\dfrac{3\cdot42}{4}-26}{10} \cdot(80-60)=71.
		\end{align*}
	}
\end{ex}

\begin{ex}%[1D1B2-1]
	Khảo sát về số giờ mượn sách thư viện của học sinh khối 11 trường Y ta được một mẫu số liệu ghép nhóm như sau:
	\begin{center}
		\begin{tabular}{|c|c|c|c|c|c|}
			\hline Số giờ mượn & {$[0 ; 4)$} & {$[4 ; 8)$} & {$[8 ; 12)$} & {$[12 ; 16)$} & {$[16 ; 20)$} \\
			\hline Số học sinh & $54$ & $78$ & $120$ & $45$ & $12$\\
			\hline
		\end{tabular}
	\end{center}
	Số trung vị của mẫu số liệu trên là
	\choice
	{\True $8,75$}
	{$8$}
	{$10$}
	{$9$}
	\loigiai{
		Ta có: Cỡ mẫu $n=309\Rightarrow\dfrac{n}{2}=154,5$.\\
		Suy ra nhóm $\left[8\,;\,12\right)$ chứa số trung vị.\\
		Do đó số trung vị của mẫu số liệu trên là\\
		$M_e=8+\dfrac{154,5-132}{120}.4=8,75$.}
\end{ex}

\begin{ex}%[1D1B2-1]
	Khảo sát về cân nặng của các học sinh lớp $11D3$ người ta được một mẫu dữ liệu ghép nhóm như sau:
	\begin{center}
		\begin{tabular}{|c|c|c|c|c|c|c|}
			\hline Cân nặng & {$[30 ; 40)$} & {$[40 ; 50)$} & {$[50 ; 60)$} & {$[60 ; 70)$} & {$[70 ; 80)$} & {$[80 ; 90)$} \\
			\hline Số học sinh & $2$ & $10$ & $16$ & $8$ & $2$ & $2$ \\
			\hline
		\end{tabular}
	\end{center}
	Số trung vị của mẫu số liệu trên là
	\choice
	{$50$}
	{$60$}
	{\True $55$}
	{$48$}
	\loigiai{
		Ta có: Số phần tử của mẫu là $n=40\Rightarrow\dfrac{n}{2}=20$.\\
		Suy ra nhóm $\left[50;60\right)$ chứa số trung vị.\\
		Do đó số trung vị của mẫu số liệu trên là\\
		$M_e=50+\dfrac{20-12}{16}\cdot10=55$.}
\end{ex}

\begin{ex}%[1D1B2-2]
	Khảo sát về cân nặng của các học sinh lớp $11$D$3$ người ta được một mẫu dữ liệu ghép nhóm như sau
	\begin{center}
		\begin{tabular}{|c|c|c|c|c|c|c|}
			\hline Cân nặng & {$[30 ; 40)$} & {$[40 ; 50)$} & {$[50 ; 60)$} & {$[60 ; 70)$} & {$[70 ; 80)$} & {$[80 ; 90)$} \\
			\hline Số học sinh & $2$ & $10$ & $16$ & $8$ & $2$ & $2$ \\
			\hline
		\end{tabular}
	\end{center}
	Tứ phân vị thứ nhất của mẫu số liệu trên là
	\choice
	{$50$}
	{\True $48$}
	{$40$}
	{$45$}
	\loigiai{
		Ta có Số phần tử của mẫu là $n=40\Rightarrow\dfrac{n}{4}=10$.\\
		Suy ra nhóm $\left[40\,;\,50\right)$ chứa tứ phân vị thứ nhất.\\
		Do đó tứ phân vị thứ nhất của mẫu số liệu trên là\\
		$Q_1=40+\dfrac{10-2}{10}\cdot10=48$.}
\end{ex}

\begin{ex}%[1D1B2-2]
	Khảo sát về cân nặng của các học sinh lớp $11$D$3$ người ta được một mẫu dữ liệu ghép nhóm như sau
	\begin{center}
		\begin{tabular}{|c|c|c|c|c|c|c|}
			\hline Cân nặng & {$[30 ; 40)$} & {$[40 ; 50)$} & {$[50 ; 60)$} & {$[60 ; 70)$} & {$[70 ; 80)$} & {$[80 ; 90)$} \\
			\hline Số học sinh & $2$ & $10$ & $16$ & $8$ & $2$ & $2$ \\
			\hline
		\end{tabular}
	\end{center}
	Tứ phân vị thứ ba của mẫu số liệu trên là
	\choice{$60$}
	{\True $62{,}5$}
	{$65$}
	{$70$}
	\loigiai{
		Ta có số phần tử của mẫu là $n=40 \Rightarrow \dfrac{3 n}{4}=30$.\\
		Suy ra nhóm $[60 ; 70)$ chứa tứ phân vị thứ ba.\\
		Do đó tứ phân vị thứ ba của mẫu số liệu trên là
		$$Q_3=60+\dfrac{30-28}{8} \cdot 10=62{,}5.$$	
	}
\end{ex}

\begin{ex}%[1D1B2-1]
	Thống kê về thời lượng mỗi trận đấu bi-a trong vòng tứ kết giải đấu European Open người ta được mẫu số liệu ghép nhóm như sau
	\begin{center}
		\begin{tabular}{|c|c|c|c|c|c|}
			\hline Thời gian & {$[9{,}5 ; 12{,}5)$} & {$[12{,}5 ; 15{,}5)$} & {$[15{,}5 ; 18{,}5)$} & {$[18{,}5 ; 21{,}5)$} & {$[21{,}5 ; 24{,}5)$} \\
			\hline Số trận & $3$ & $12$ & $15$ & $24$ & $2$ \\
			\hline
		\end{tabular}
	\end{center}
	Số trung vị của mẫu số liệu trên là
	\choice{$18$}
	{\True $18{,}1$}
	{$18{,}2$}
	{$18{,}3$}
	\loigiai{
		Ta có số phần tử của mẫu là $n=56 \Rightarrow \dfrac{n}{2}=28$.\\
		Suy ra nhóm $[15{,}5 ; 18{,}5)$ chứa số trung vị.\\
		Do đó số trung vị của mẫu số liệu trên là
		$M_e=15{,}5+\dfrac{28-15}{15} \cdot 3=18{,}1$. 
	}
\end{ex}

\begin{ex}%[1D1B2-2]
	Thống kê về thời lượng mỗi trận đấu bi-a trong vòng tứ kết giải đấu European Open người ta được mẫu số liệu ghép nhóm như sau:
	\begin{center}
		\begin{tabular}{|c|c|c|c|c|c|}
			\hline Thời gian & {$[9{,}5 ; 12{,}5)$} & {$[12{,}5 ; 15{,}5)$} & {$[15{,}5 ; 18{,}5)$} & {$[18{,}5 ; 21{,}5)$} & {$[21{,}5 ; 24{,}5)$} \\
			\hline Số trận & $3$ & $12$ & $15$ & $24$ & $2$ \\
			\hline
		\end{tabular}
	\end{center}
	Tứ phân vị thứ ba của mẫu số liệu trên là
	\choice{$19$}
	{$19{,}5$}		
	{\True $20$}
	{$20{,}5$}
	\loigiai{
		Ta có số phần tử của mẫu là $n=56 \Rightarrow \dfrac{3 n}{4}=42$.\\
		Suy ra nhóm $[18{,}5 ; 21{,}5)$ chứa tứ phân vị thứ ba.\\
		Do đó tứ phân vị thứ ba của mẫu số liệu trên là
		$$
		Q_3=18{,}5+\dfrac{42-30}{24} \cdot 3=20 .
		$$
	}
\end{ex}

\begin{ex}%[1D1B2-1]
	Doanh thu bán hàng trong $20$ ngày được lựa chọn ngẫu nhiên của một của hàng được ghi lại ở bảng sau (đơn vị: triệu đồng)
	\begin{center}
		\begin{tabular}{|c|c|c|c|c|c|}
			\hline Doanh thu & {$[5 ; 7)$} & {$[7 ; 9)$} & {$[9 ; 11)$} & {$[11 ; 13)$} & {$[13 ; 15)$} \\
			\hline Số ngày & $2$ & $7$ & $7$ & $3$ & $1$ \\
			\hline
		\end{tabular}
	\end{center}
	Trung vị của mẫu số liệu trên là?
	\choice{$6{,}286$}
	{ \True $9{,}286$}
	{ $9{,}386$}
	{ $6{,}386$}
	\loigiai{
		Gọi $x_1$, $x_2$, $\ldots$, $x_{20}$ là doanh thu bán hàng trong 20 ngày xếp theo thứ tự không giảm.\\
		Khi đó $x_1$, $x_2 \in[5 ; 7)$; $x_3, \ldots, x_9 \in[7 ; 9)$; $x_9$, $\ldots$, $x_{16} \in[9 ; 11)$; $x_{17}$, $\ldots$, $x_{19} \in[11 ; 13)$; $x_{20} \in[13 ; 15)$.\\
		Do đó trung vị của mẫu số liệu thuộc nhóm $[9 ; 11)$.\\
		Ta xác định được $n=20$, $n_m=7$, $C=2+7=9$, $U_m=9$, $U_{m+1}=11$.\\
		Vậy trung vị của mẫu ghép nhóm là
		$$
		M_e=9+\dfrac{\dfrac{20}{2}-9}{7} \cdot (11-9)=9{,}286.
		$$
	}
\end{ex}


\begin{ex}%[1D1B2-2]
	Doanh thu bán hàng trong $20$ ngày được lựa chọn ngẫu nhiên của một của hàng được ghi lại ở bảng sau (đơn vị: triệu đồng)
	\begin{center}
		\begin{tabular}{|c|c|c|c|c|c|}
			\hline Doanh thu & {$[5 ; 7)$} & {$[7 ; 9)$} & {$[9 ; 11)$} & {$[11 ; 13)$} & {$[13 ; 15)$} \\
			\hline Số ngày & $2$ & $7$ & $7$ & $3$ & $1$ \\
			\hline
		\end{tabular}
	\end{center}
	Tứ phân vị thứ nhất của mẫu số liệu gần nhất với giá trị nào trong các giá trị dưới đây?
	\choice{ $7$}
	{$7{,}6$}
	{ \True $8$}
	{ $8{,}6$}
	\loigiai{
		Gọi $x_1$, $x_2$, $\ldots$, $x_{20}$ là doanh thu bán hàng trong 20 ngày xếp theo thứ tự không giảm.\\
		Khi đó $x_1$, $x_2 \in[5 ; 7)$; $x_3, \ldots, x_9 \in[7 ; 9)$; $x_9$, $\ldots$, $x_{16} \in[9 ; 11)$; $x_{17}$, $\ldots$, $x_{19} \in[11 ; 13)$; $x_{20} \in[13 ; 15)$.\\
		Do đó, tứ phân vị thứ nhất của mẫu số liệu thuộc nhóm $[7 ; 9)$.\\
		$n=20$, $n_m=7$, $C=2$, $u_m=7$, $u_{m+1}=9$.
		$$
		Q_1=7+\dfrac{\dfrac{1.20}{4}-2}{7}(9-7) \approx 7{,}86 \approx 8.
		$$
		
	}
\end{ex}


\begin{ex}%[1D1B2-2]
	Doanh thu bán hàng trong $20$ ngày được lựa chọn ngẫu nhiên của một của hàng được ghi lại ở bảng sau (đơn vị: triệu đồng)
	\begin{center}
		\begin{tabular}{|c|c|c|c|c|c|}
			\hline Doanh thu & {$[5 ; 7)$} & {$[7 ; 9)$} & {$[9 ; 11)$} & {$[11 ; 13)$} & {$[13 ; 15)$} \\
			\hline Số ngày & $2$ & $7$ & $7$ & $3$ & $1$ \\
			\hline
		\end{tabular}
	\end{center}
	Tứ phân vị thứ ba của mẫu số liệu gần nhất với giá trị nào trong các giá trị dưới đây?
	\choice{ $10$}
	{\True $11$}
	{  $12$}
	{ $13$}
	\loigiai{
		Gọi $x_1$, $x_2$, $\ldots$, $x_{20}$ là doanh thu bán hàng trong 20 ngày xếp theo thứ tự không giảm.\\
		Khi đó $x_1$, $x_2 \in[5 ; 7)$; $x_3, \ldots, x_9 \in[7 ; 9)$; $x_9$, $\ldots$, $x_{16} \in[9 ; 11)$; $x_{17}$, $\ldots$, $x_{19} \in[11 ; 13)$; $x_{20} \in[13 ; 15)$.\\
		Do đó, tứ phân vị thứ ba của mẫu số liệu thuộc nhóm $[9 ; 11)$.\\
		$
		n=20$, $n_m=7$, $C=9$, $u_m=9$, $u_{m+1}=11
		$.
		$$
		Q_3=9+\dfrac{\dfrac{3\cdot20}{4}-9}{7}(11-9) \approx 10{,}71 \approx 11.
		$$
	}
\end{ex}

\begin{ex}%[1D1K2-1]
	Trung tâm ngoại ngữ thống kê bảng điểm môn Tiếng Anh của một khóa học trong bảng bên dưới
	\begin{center}
		\begin{tabular}{|l|c|c|c|c|c|}
			\hline
			Điểm     & [0;2) & [2;4) & [4;6) & [6;8) & [8;10) \\ \hline
			Học viên & 10   & 30   & 55   & 42   & 9     \\ \hline
		\end{tabular}
	\end{center}
	Trung vị của mẫu số liệu thuộc khoảng nào trong các khoảng dưới đây?
	\choice
	{$[2;4)$}
	{$[8;10)$}
	{\True $[4;6)$}
	{$[0;2)$}
	\loigiai{
		Gọi $x_{1}, x_{2}, ..., x_{146}$ là số liệu được sắp xếp theo thứ tự không giảm. \\
		Ta có:
		$x_{1},...,x_{10}\in [0;2); x_{11},...,x_{40}\in [2;4); x_{41},...,x_{95}\in [4;6); x_{96},...,x_{137}\in [6;8); x_{138},...,x_{146}\in [8;10)$
		nên trung vị của mẫu số liệu $x_{1},x_{2},...,x_{146}$ là $\dfrac{1}{2}(x_{73}+x_{74})\in [4;6)$. \\
		Suy ra $[4;6)$ là nhóm chứa trung vị. Khi đó ta xác định được \\
		$n=146,n_{m}=55,C=10+30=40,u_{m}=4,u_{m+1}=6$ \\
		Vậy trung vị vủa mẫu số liệu ghép nhóm là \\		
		$$M_{e}=u_{m}+\dfrac{\dfrac{n}{2}-C}{n_{m}}.(u_{m+1}-u_{m})=4+\dfrac{\dfrac{146}{2}-40}{55}.(6-4)=5,2.$$
	}
\end{ex}

\begin{ex}%[1D1B2-2]
	Trung tâm ngoại ngữ thống kê bảng điểm môn Tiếng Anh của một khóa học trong bảng bên dưới
	\begin{center}
		\begin{tabular}{|l|c|c|c|c|c|}
			\hline
			Điểm     & [0;2) & [2;4) & [4;6) & [6;8) & [8;10) \\ \hline
			Học viên & 10   & 30   & 55   & 42   & 9     \\ \hline
		\end{tabular}
	\end{center}
	Tứ phân vị thứ nhất của mẫu số liệu trên gần nhất với giá trị nào trong các giá trị sau?
	\choice
	{\True $3,8$}
	{$3,5$}
	{$2,2$}
	{$3,9$}
	\loigiai{
		Gọi $x_{1}, x_{2}, ..., x_{146}$ là số liệu được sắp xếp theo thứ tự không giảm. \\
		Tứ phân vị thứ nhất của của dãy số liệu $x_{1},x_{2},...x_{146}$ là $x_{37}\in [2;4)$. Do đó, tứ phân vị thứ nhất của mẫu số liệu ghép nhóm trên là \\
		$$Q_{1}=2+\dfrac{\dfrac{1.146}{4}-10}{30}.(4-2)\approx 3,8.$$
	}
\end{ex}
\begin{ex}%[1D1B2-2]
	Trung tâm ngoại ngữ thống kê bảng điểm môn Tiếng Anh của một khóa học trong bảng bên dưới
	\begin{center}
		\begin{tabular}{|l|c|c|c|c|c|}
			\hline
			Điểm     & [0;2) & [2;4) & [4;6) & [6;8) & [8;10) \\ \hline
			Học viên & 10   & 30   & 55   & 42   & 9     \\ \hline
		\end{tabular}
	\end{center}
	Tứ phân vị thứ ba của mẫu số liệu trên gần nhất với giá trị nào trong các giá trị sau?
	\choice
	{$6,1$}
	{\True $6,7$}
	{$6,9$}
	{$6,5$}
	\loigiai{
		Gọi $x_{1}, x_{2}, ..., x_{146}$ là số liệu được sắp xếp theo thứ tự không giảm. \\
		Tứ phân vị thứ nhất của của dãy số liệu $x_{1},x_{2},...x_{146}$ là $x_{110}\in [6;8)$. Do đó, tứ phân vị thứ nhất của mẫu số liệu ghép nhóm trên là \\
		$$Q_{3}=6+\dfrac{\dfrac{3.146}{4}-(10+30+35)}{42}.(8-6)\approx 6,7.$$
	}
\end{ex}

\begin{ex}%[1D1B2-1]
	\sochc{3}
	Kết quả điều tra chiều cao học sinh (đơn vị: cm) trong trường THPT được biểu diễn ở biểu đồ dưới đây.
	\begin{center}
	\begin{tikzpicture}[line join=round, line cap=round, thick, >=stealth, scale=0.5]
		\draw[->] (0,-1)--(0,14)node[left]{(Số học sinh)};
		\draw (-1,0)--(0,0)node[below left]{$O$}--(3,0); 
		\draw[dotted] (3,0)--(4,0);
		\draw[->] (4,0)--(22,0)node[below]{(cm)};
		\foreach \a in {1,...,20}{\draw[line width=1pt,shift={(\a,0)}] (0pt,2pt)--(0pt,-2pt);} %Chia đơn vị trục Ox
		\foreach \i/\y/\g/\h in 
		{0/6/horizontal lines/[160;163),
			1/12/grid/[163;166),
			2/9/vertical lines/[166;169),
			3/5/dots/[169;172),
			4/3/bricks/[172;175)}
		{\draw[line width=1.5pt, pattern=\g](5+3*\i,0)rectangle++(3,\y); %Vẽ cột
			\draw[dashed] (3*\i+5,\y)--(0,\y); 	%Vẽ nét đứt gióng cột
			\draw[line width=1pt,shift={(0,\y)}] (2pt,0pt)--(-2pt,0pt); %Chia đơn vị Oy
			\draw(0,\y) node[left] {$\y$}; %Kí hiệu giá trị trục Oy
			\draw(6.5+3*\i,0) node[below] {\h};} %Kí hiệu nhãn cột
		\path (current bounding box.north) node[]{Chiều cao của 35 học sinh trường THPT};
	\end{tikzpicture}
\end{center}
\begin{chc}
	Trung vị của mẫu số liệu ghép nhóm trên là bao nhiêu (Đáp án làm tròn đến hàng phần trăm).
	\choice
	{161,88}
	{167,85}
	{1663,87}
	{\True165,88}
	\loigiai{ 
		Gọi $x_1;x_2;...;x_{35}$ là cân nặng của các học sinh trường THPT xếp theo thứ tự không giảm.\\ 
		Do $x_1;...;x_6\in[160;163);x_{7};...;x_{18}\in[163;166);x_{19};...;x_{27}\in[166;169);x_{28};...;x_{32}\in[169;172);\\x_{33};_{34};x_{35}\in[172;175)$ nên trung vị của mẫu số liệu $x_1;x_2;...;x_{35}$ là $x_{18}\in[163;166)$.\\
		Ta xác định được $n=35,n_m=12,C=6,u_m=163;u_{m+1}=166$.\\
		Vậy trung vị của mẫu số liệu ghép nhóm là
		$$M_e=163+\dfrac{\dfrac{35}{2}-6}{12}.(166-163)\approx165,88.$$}
\end{chc}
\begin{chc}
	Tứ phân vị thứ nhất của mẫu số liệu ghép nhóm trên là bao nhiêu? (Đáp án làm tròn đến hàng phần trăm).
	\choice
	{161,69}
	{167,69}
	{165,69}
	{\True163,69}
	\loigiai{ Gọi $x_1;x_2;...;x_{35}$ là cân nặng của các học sinh trường THPT xếp theo thứ tự không giảm.\\ Do 
		$x_1;...;x_6\in[160;163);x_{7};...;x_{18}\in[163;166);x_{19};...;x_{27}\in[166;169);x_{28};...;x_{32}\in[169;172);\\x_{33};_{34};x_{35}\in[172;175)$ nên tứ phân vị thứ nhất của dãy số liệu $x_1;x_2;...;x_{35}$ là $x_{9}\in[163;166)$.\\
		Vậy tứ phân vị thứ nhất của mẫu số liệu ghép nhóm là
		$$Q_1=163+\dfrac{\dfrac{35}{4}-6}{12}.(166-163)\approx163,688.$$}
\end{chc}
\begin{chc}
	Tứ phân vị thứ ba của mẫu số liệu ghép nhóm trên là bao nhiêu?
	\choice
	{162,75}
	{164,75}
	{166,75}
	{\True168,75}
	\loigiai{ Gọi $x_1;x_2;...;x_{35}$ là cân nặng của các học sinh trường THPT xếp theo thứ tự không giảm.\\ Do 
		$x_1;...;x_6\in[160;163);x_{7};...;x_{18}\in[163;166);x_{19};...;x_{27}\in[166;169);x_{28};...;x_{32}\in[169;172);\\x_{33};_{34};x_{35}\in[172;175)$ nên tứ phân vị thứ ba của dãy số liệu $x_1;x_2;...;x_{35}$ là $x_{27}\in[166;169)$.\\
		Vậy tứ phân vị thứ ba của mẫu số liệu ghép nhóm là
		$$Q_3=166+\dfrac{\dfrac{3.35}{4}-(6+12)}{9}.(169-166)=168,75.$$}
\end{chc}
\end{ex}
% \begin{ex}%[1D1B2-2]
% 	Kết quả điều tra chiều cao học sinh (đơn vị: cm) trong trường THPT được biểu diễn ở biểu đồ dưới đây.
% 	\begin{center}
% 	\begin{tikzpicture}[line join=round, line cap=round, thick, scale=0.5]
% 		\draw[->] (0,-1)--(0,14)node[left]{(Số học sinh)};
% 		\draw (-1,0)--(0,0)node[below left]{$O$}--(3,0); 
% 		\draw[dotted] (3,0)--(4,0);
% 		\draw[->] (4,0)--(22,0)node[below]{(cm)};
% 		\foreach \a in {1,...,20}{\draw[line width=1pt,shift={(\a,0)}] (0pt,2pt)--(0pt,-2pt);} %Chia đơn vị trục Ox
% 		\foreach \i/\y/\g/\h in 
% 		{0/6/horizontal lines/[160;163),
% 			1/12/grid/[163;166),
% 			2/9/vertical lines/[166;169),
% 			3/5/dots/[169;172),
% 			4/3/bricks/[172;175)}
% 		{\draw[line width=1.5pt, pattern=\g](5+3*\i,0)rectangle++(3,\y); %Vẽ cột
% 			\draw[dashed] (3*\i+5,\y)--(0,\y); 	%Vẽ nét đứt gióng cột
% 			\draw[line width=1pt,shift={(0,\y)}] (2pt,0pt)--(-2pt,0pt); %Chia đơn vị Oy
% 			\draw(0,\y) node[left] {$\y$}; %Kí hiệu giá trị trục Oy
% 			\draw(6.5+3*\i,0) node[below] {\h};} %Kí hiệu nhãn cột
% 		\path (current bounding box.north) node[]{Chiều cao của 35 học sinh trường THPT};
% 	\end{tikzpicture}
% \end{center}
% 	Tứ phân vị thứ nhất của mẫu số liệu ghép nhóm trên là bao nhiêu? (Đáp án làm tròn đến hàng phần trăm).
% 	\choice
% 	{161,69}
% 	{167,69}
% 	{165,69}
% 	{\True163,69}
% 	\loigiai{ Gọi $x_1;x_2;...;x_{35}$ là cân nặng của các học sinh trường THPT xếp theo thứ tự không giảm.\\ Do 
% 		$x_1;...;x_6\in[160;163);x_{7};...;x_{18}\in[163;166);x_{19};...;x_{27}\in[166;169);x_{28};...;x_{32}\in[169;172);\\x_{33};_{34};x_{35}\in[172;175)$ nên tứ phân vị thứ nhất của dãy số liệu $x_1;x_2;...;x_{35}$ là $x_{9}\in[163;166)$.\\
% 		Vậy tứ phân vị thứ nhất của mẫu số liệu ghép nhóm là
% 		$$Q_1=163+\dfrac{\dfrac{35}{4}-6}{12}.(166-163)\approx163,688.$$}
% \end{ex}
% \begin{ex}%[1D1B2-2]
% 	Kết quả điều tra chiều cao học sinh (đơn vị: cm) trong trường THPT được biểu diễn ở biểu đồ dưới đây.
% 	$$	\begin{tikzpicture}[line join=round, line cap=round, thick, scale=0.6]
% 		\draw[->] (0,-1)--(0,14)node[left]{(Số học sinh)};
% 		\draw (-1,0)--(0,0)node[below left]{$O$}--(3,0); 
% 		\draw[dotted] (3,0)--(4,0);
% 		\draw[->] (4,0)--(22,0)node[below]{(cm)};
% 		\foreach \a in {1,...,20}{\draw[line width=1pt,shift={(\a,0)}] (0pt,2pt)--(0pt,-2pt);} %Chia đơn vị trục Ox
% 		\foreach \i/\y/\g/\h in 
% 		{0/6/horizontal lines/[160;163),
% 			1/12/grid/[163;166),
% 			2/9/vertical lines/[166;169),
% 			3/5/dots/[169;172),
% 			4/3/bricks/[172;175)}
% 		{\draw[line width=1.5pt, pattern=\g](5+3*\i,0)rectangle++(3,\y); %Vẽ cột
% 			\draw[dashed] (3*\i+5,\y)--(0,\y); 	%Vẽ nét đứt gióng cột
% 			\draw[line width=1pt,shift={(0,\y)}] (2pt,0pt)--(-2pt,0pt); %Chia đơn vị Oy
% 			\draw(0,\y) node[left] {$\y$}; %Kí hiệu giá trị trục Oy
% 			\draw(6.5+3*\i,0) node[below] {\h};} %Kí hiệu nhãn cột
% 		\path (current bounding box.north) node[]{Chiều cao của 35 học sinh trường THPT};
% 	\end{tikzpicture}$$
% 	Tứ phân vị thứ ba của mẫu số liệu ghép nhóm trên là bao nhiêu?
% 	\choice
% 	{162,75}
% 	{164,75}
% 	{166,75}
% 	{\True168,75}
% 	\loigiai{ Gọi $x_1;x_2;...;x_{35}$ là cân nặng của các học sinh trường THPT xếp theo thứ tự không giảm.\\ Do 
% 		$x_1;...;x_6\in[160;163);x_{7};...;x_{18}\in[163;166);x_{19};...;x_{27}\in[166;169);x_{28};...;x_{32}\in[169;172);\\x_{33};_{34};x_{35}\in[172;175)$ nên tứ phân vị thứ ba của dãy số liệu $x_1;x_2;...;x_{35}$ là $x_{27}\in[166;169)$.\\
% 		Vậy tứ phân vị thứ ba của mẫu số liệu ghép nhóm là
% 		$$Q_3=166+\dfrac{\dfrac{3.35}{4}-(6+12)}{9}.(169-166)=168,75.$$}
% \end{ex}
\Closesolutionfile{ans}

\ind{PHẦN II.} \inden{Câu trắc nghiệm đúng sai. Trong mỗi ý a), b), c), d) ở mỗi câu, học sinh chọn đúng hoặc sai.}\\
\Opensolutionfile{ans}[ans/1D3-B2-2]
\begin{ex}%[1D3H2-3]
	Cho bảng số liệu sau
	\begin{center}
		\begin{tabular}{|c|c|c|c|c|c|c|}
			\hline
			Nhóm & $[30;40)$ & $[40;50)$ & $[50;60)$ & $[60;70)$ & $[70;80)$ & $[80;90)$ \\
			\hline
			Tần số & $2$ & $10$ & $16$ & $8$ & $2$ & $2$ \\
			\hline
		\end{tabular}
	\end{center}
	\choiceTF
	{\True Cỡ mẫu của mẫu số liệu là $n=40$}
	{\True Tứ phân vị thứ nhất của mẫu số liệu ghép nhóm là $Q_1=48$}
	{Tứ phân vị thứ hai của mẫu số liệu ghép nhóm là $Q_2=45$}
	{Tứ phân vị thứ ba của mẫu số liệu ghép nhóm là $Q_3=61{,}5$}
	\loigiai{
		\begin{itemchoice}
			\itemch Cỡ mẫu của mẫu số liệu là $n=40$.
			\itemch Xét nửa mẫu số liệu bên trái $x_1, x_2, x_3, \ldots, x_{20}$ có trung vị $\dfrac{x_{10}+x_{11}}{2} \in [40;50)$.\\
			Ta có $n_i=10$; $C_1=2$; $u_i=40$; $u_{i+1}=50$.\\
			Tứ phân vị thứ nhất của mẫu số liệu ghép nhóm là
			\[
			Q_1=40+\dfrac{\dfrac{40}{4}-2}{10} (50-40)=48.
			\]
			\itemch Gọi $x_1, x_2, x_3, \ldots, x_{40}$ là mẫu số liệu được xếp theo thứ tự không giảm.\\
			Trung vị của mẫu số liệu $\dfrac{x_{20}+x_{21}}{2} \in [50;60)$.\\
			Ta có $n_m=16$; $C_2=2+10=12$; $u_m=50$; $u_{n+1}=60$.\\
			Tứ phân vị thứ hai của mẫu số liệu ghép nhóm cũng là trung vị của mẫu số liệu đó là
			\[
			Q_2=M_e=50+\dfrac{\dfrac{40}{2}-12}{16} (60-50)=55.
			\]
			\itemch Xét nửa mẫu số liệu bên phải $x_{21}, x_{22}, x_{23}, \ldots, x_{40}$ có trung vị $\dfrac{x_{30}+x_{31}}{2} \in [60;70)$.\\
			Ta có $n_j=8$; $C_3=2+10+16=28$; $u_j=60$; $u_{j+1}=70$.\\
			Tứ phân vị thứ ba của mẫu số liệu ghép nhóm là
			\[
			Q_3=60+\dfrac{\dfrac{3\cdot 40}{4}-28}{8} (70-60)=62{,}5.
			\]
		\end{itemchoice}
	}
\end{ex}

% \begin{ex}
% 	Cho mẫu số liệu ghép nhóm về quãng đường chạy bộ(đơn vị: km) và số ngày chạy bộ như hình dưới đây.Xét tính đúng-sai của các khẳng định sau (kết quả làm tròn đến hàng phần trăm).
% 	\begin{center}\begin{tabular}{|c|c|c|c|c|c|c|}
% 			\hline
% 			Quãng đường chạy bộ(đơn vị: km)   & [1,2 ; 1,7) & [1,7 ; 2,2) & [2,2 ; 2,7) & [2,7 ; 3,2) & [3,2 ; 3,7) & [3,7 ; 4,2)\\  
% 			\hline 
% 			Số ngày chạy bộ & 2 & 3 & 6 & 7 & 5 & 6 \\ 
% 			\hline 
% 		\end{tabular}
% 	\end{center}
% 	\choiceTF
% 	{  Giá trị ${3,57}$ là một mốt của mẫu số liệu }
% 	{ \True  Số trung bình của mẫu số liệu ghép nhóm là ${2,93}$ }
% 	{ \True  Tứ phân vị $Q_1$ của mẫu số liệu là ${2,39}$ }
% 	{ \True  Tứ phân vị $Q_3$ của mẫu số liệu là ${3,58}$ }
% 	\loigiai{ 
		
		
% 		a) Khẳng định đã cho là khẳng định sai.
		
% 		Nhóm chứa mốt là nhóm thứ $k=4, n_k=7,n_{k-1}=6, n_{k+1}=5, u_k=2.7, u_{k+1}=3.2$.
		
% 		Áp dụng công thức tìm Mốt: $M_o=u_k+\dfrac{ (n_k - n_{k-1}) }{(n_k-n_{k-1})+(n_k-n_{k+1}) }(u_{k+1}-u_k)=2,87$.
		
% 		b) Khẳng định đã cho là khẳng định đúng.
		
% 		Các giá trị đại diện của mẫu số liệu là: 1,45; 1,9500000000000002; 2,45; 2,95; 3,45; 3,95
		
% 		Tổng tần số là: $n=29$
		
% 		Số trung bình của mẫu số liệu ghép nhóm là:
		
% 		$\overline{x}=\dfrac{2.1,45+3.1,9500000000000002+6.2,45+7.2,95+5.3,45+6.3,95}{ 29}=2,93$.
		
		
		
% 		c) Khẳng định đã cho là khẳng định đúng.
		
% 		Bước 1: Tổng tần số là: $N=29$.
		
% 		Bước 2: Xác định vị trí của $Q_1$: $Q_1$ nằm ở vị trí $\dfrac{29}{4}=7.3$.
		
% 		Bước 3: Xác định lớp chứa $Q_1$: tính tần số tích lũy từ lớp đầu tiên đến khi đạt hoặc vượt qua vị trí của $Q_1$ ta được lớp $[2.2;2.7)$.
		
% 		Bước 4: Xác định các thông số của công thức tính $Q_1$.
		
% 		Cận dưới của lớp chứa $Q_1$: $L=2.2$
		
% 		Tổng tần số của các lớp trước lớp chứa $Q_1$: $F=5$
		
% 		Tần số của lớp chứa $Q_1$: $f=6$.
		
% 		Độ rộng lớp chứa $Q_1$: $h=2.7 - 2.2=0.5$.
		
% 		Áp dụng công thức: $Q_1=L+\left(\dfrac{ \dfrac{N}{4}-F }{f}\right).h=2.2+\left(\dfrac{ \dfrac{29}{4}-5 }{6}\right).0.5=2,39$.
		
% 		d) Khẳng định đã cho là khẳng định đúng.
		
% 		Bước 1: Tổng tần số là: $N=29$.
		
% 		Bước 2: Xác định vị trí của $Q_3$: $Q_3$ nằm ở vị trí $\dfrac{3.29}{4}=21.8$.
		
% 		Bước 3: Xác định lớp chứa $Q_3$: tính tần số tích lũy từ lớp đầu tiên đến khi đạt hoặc vượt qua vị trí của $Q_3$ ta được lớp $[3.2;3.7)$.
		
% 		Bước 4: Xác định các thông số của công thức tính $Q_3$.
		
% 		Cận dưới của lớp chứa $Q_3$: $L=3.2$
		
% 		Tổng tần số của các lớp trước lớp chứa $Q_3$: $F=18$
		
% 		Tần số của lớp chứa $Q_3$: $f=5$.
		
% 		Độ rộng lớp chứa $Q_3$: $h=3.7 - 3.2=0.5$.
		
% 		Áp dụng công thức: $Q_3=L+\left(\dfrac{ \dfrac{3N}{4}-F }{f}\right).h=3.2+\left(\dfrac{ \dfrac{3.29}{4}-18 }{5}\right).0.5=3,58$.
		
		
% }\end{ex}

\begin{ex}%[1D3V1-2]%[1D3V2-3]%[1-TN-DS-TLN-THPT-NguyenThiMinhKhai-HCM-HKI-NH24-25]%[Đức Dương]
	Bảng thống kê doanh số bán hàng của các nhân viên một siêu thị điện máy trong một dịp Black Friday như sau:
	\begin{center}
		\begin{tabular}{|c|c|c|c|c|c|}
			\hline \begin{tabular}{c} 
				Doanh số \\
				(triệu đồng)
			\end{tabular} & $[20;30)$ & $[30;40)$ & $[40;50)$ & $[50;60)$ & $[60;70)$ \\
			\hline Số nhân viên & $4$ & $8$ & $12$ & $7$ & $5$ \\
			\hline
		\end{tabular}
	\end{center}
	\choiceTF
	{Cỡ mẫu của mẫu số liệu trên là $n=35$}
	{\True Số trung bình của mẫu số liệu trên là $\overline{x}\approx 45{,}28$ triệu đồng (\textit{kết quả làm tròn đến hàng phần trăm})}
	{Tứ phân vị thứ nhất của mẫu số liệu trên là $Q_1=36{,}5$}
	{Cửa hàng dự định sẽ thưởng $25\%$ số nhân viên có doanh số bán hàng cao nhất. Theo mẫu sổ liệu trên, của hàng nên khen thưởng các nhân viên có doanh số bán hàng ít nhất từ $53{,}29$ triệu đồng (\textit{kết quả làm tròn đến hàng phần trăm})}
	\loigiai{
		Theo bài ra ta có bảng tần số ghép nhóm sau
		\begin{center}
			\begin{tabular}{|c|c|c|c|c|c|}
				\hline \begin{tabular}{c} 
					Doanh số \\
					(triệu đồng)
				\end{tabular} & $[20;30)$ & $[30;40)$ & $[40;50)$ & $[50;60)$ & $[60;70)$ \\
				\hline Giá trị đại diện & $25$ & $35$ & $45$ & $55$ & $65$ \\		
				\hline Số nhân viên & $4$ & $8$ & $12$ & $7$ & $5$ \\
				\hline
			\end{tabular}
		\end{center}
		\begin{itemchoice}
			\itemch \textbf{Sai}.\\
			Cỡ mẫu của mẫu số liệu trên là $n=4+8+12+7+5=36$.
			\itemch \textbf{Đúng}.\\
			Giá trị trung bình của mẫu số liệu 
			\[ \overline{x}=\dfrac{25\cdot 4+35\cdot 8+45\cdot 12+55\cdot 7+65\cdot 5}{36}=\dfrac{815}{18}\approx 45{,}28.\]
			\itemch \textbf{Sai}.\\
			Ta có $\dfrac{n}{4}=9$ mà $4<9<12$ nên suy ra tứ phân vị thứ nhất là 
			\[ Q_1=30+\dfrac{9-4}{8}\cdot 10=36{,}25.\]
			\itemch \textbf{Sai}.\\
			Ta có $\dfrac{3n}{4}=27$ mà $4+8+12<27<4+8+12+7$ nên suy ra tứ phân vị thứ ba là 
			\[ Q_3=50+\dfrac{27-24}{7}\cdot 10=\dfrac{380}{7}\approx 54{,}29. \]
			Vậy cửa hàng nên khen thưởng các nhân viên có doanh số bán hàng ít nhất là $54{,}29$ triệu đồng.
		\end{itemchoice}
	}
\end{ex}

\begin{ex}%[1D3V2-3]%[TH-THCS-THPT-HoangViet-DakLak-HKI-NH24-25]%[Trần Hưng]
	Một bảng xếp hạng đã tính điểm chuẩn hoá cho chỉ số nghiên cứu của một số trường đại học ở Việt Nam và thu được kết quả sau:
	\begin{center}
		\begin{tabular}{|c|c|c|c|c|c|c|}
			\hline
			Điểm & $[10;20)$ & $[20;30)$ & $[30;40)$ & $[40;50)$ & $[50,60)$ & $[60,70)$\\
			\hline
			Số trường & $4$ & $19$ & $6$ & $2$ & $3$ & $1$\\
			\hline
		\end{tabular}
	\end{center}
	\choiceTF
	{\True Số liệu đã cho cho có $35$ mẫu số liệu}
	{Số trung bình của mẫu số liệu đã cho là $28$}
	{Số trung vị của mẫu số liệu là $M_{e}=12$}
	{\True Ngưỡng điểm để đưa ra danh sách $25\%$ trường đại học có chỉ số nghiên cứu tốt nhất Việt Nam là trên $35{,}42$}
	\loigiai{
		\begin{itemchoice}
			\itemch 
			Số các số liệu đã cho có mẫu số liệu là $n=4+19+6+2+3+1=35$.
			\itemch 
			Số trung bình cộng của mẫu số liệu đã cho là\\
			\[\overline{x}=\dfrac{4\cdot 15+19\cdot 25+6\cdot 35+2\cdot 45+3\cdot 55+1\cdot 65}{35}\approx30{,}43.\]
			\itemch 
			Vị trí của số trung vị là $\dfrac{35+1}{2}=18$.\\
			Do đó, giá trị trung vị của mẫu số liệu nằm trong khoảng $[20;30)$.
			\itemch 
			Gọi $x_1$, $x_2$, $\ldots$, $x_{35}$ là điểm chuẩn hóa cho chỉ số nghiên cứu của các trường đại học và giả sử dãy này đã sắp xếp theo thứ tự tăng dần. Khi đó, trung vị của mẫu số liệu là $M_e=x_{18}$ và tứ phân vị thứ ba $Q_3$ của mẫu số liệu là trung vị của nửa số liệu bên phải $M_e$, đó là dãy gồm 17 số liệu $x_{19}$, $x_{20}$, $\ldots$, $x_{35}$, do đó $Q_3=x_{27}$.\\
			Do $x_{27}$ thuộc nhóm $[30;40$ ) nên nhóm này chứa $Q_3$.\\
			Do đó, $p=3$; $a_3=30$; $m_3=6$; $m_1+m_2=4+19=23$; $a_4-a_3=40-30=10$ và ta có 		
			\[Q_3=30+\dfrac{\dfrac{3{,}35}{4}-23}{6} \cdot 10 \approx 35{,}42.\]	
			$25\%$ trường đại học có chỉ số nghiên cứu tốt nhất là $25\%$ giá trị nằm bên phải $Q_3$.\\
			Vậy điểm ngưỡng để đưa ra danh sách $25\%$ trường đại học có chỉ số nghiên cứu tốt nhất Việt Nam là những trường có điểm chuẩn hóa lớn hơn hoặc bằng $35{,}42$.
		\end{itemchoice}
	}
\end{ex}

\Closesolutionfile{ans}

\ind{PHẦN III.} \inden{Câu trắc nghiệm trả lời ngắn.}\\
\Opensolutionfile{ans}[ans/1D3-B2-3]
\begin{ex}%[1D3H2-2]
	Kết quả thu thập điểm số môn Toán của $25$ học sinh khi tham gia kì thi học sinh giỏi toán $11$ (thang điểm $20$) cho ta bảng tần số ghép nhóm sau
	\begin{center}
		\begin{tabular}{|c|c|c|c|c|c|}
			\hline
			Nhóm & $[0;4)$ & $[4;8)$ & $[8;12)$ & $[12;16)$ & $[16;20)$ \\
			\hline
			Số học sinh  & $1$ & $7$ & $12$ & $3$ & $2$ \\
			\hline
		\end{tabular}
	\end{center}
	Trung vị của mẫu số liệu ghép nhóm trên bằng bao nhiêu?\\
	\shortans[oly]{$9{,}5$}
	\loigiai{
		Cỡ mẫu của mẫu số liệu là $n=25$.\\
		Gọi $x_1, x_2, x_3, \ldots, x_{25}$ là điểm số của $25$ học sinh trong kì thi đó và dãy này được sắp xếp theo thứ tự không giảm.\\
		Trung vị của mẫu số liệu là $x_{13} \in [8;12)$.\\
		Ta có $n=25$, $n_m=12$, $C=1+7=8$, $u_m=8$, $u_{m+1}=12$.\\ 
		Trung vị mẫu số liệu ghép nhóm là
		\[
		M_e=u_m+\dfrac{\dfrac{n}{2}-C}{n_m} \left(u_{m+1}-u_m \right)=8+\dfrac{\dfrac{25}{2}-8}{12} (12-8)=9{,}5.
		\]
	}
\end{ex}

\begin{ex}%[1D3H2-3][Dự án đề kiểm tra môn Toán khối 11 HKI NH24-25 - Đợt 5 - Hieu Hieu Minh Minh]%[THPT Chuyên Lê Hồng Phong - Tp HCM]
	Thời gian (phút) truy cập Internet mỗi buổi tối của một số học sinh được cho trong bảng sau:
	\begin{longtable}{|c|c|c|c|c|c|}
		\hline
		Thời gian (giờ)&$[9{,}5; 12{,}5)$ & $[12{,}5; 15{,}5)$ & $[15{,}5;18{,}5)$ & $[18{,}5;21{,}5)$ &$[21{,}5;24{,}5)$  \\
		\hline
		Số học sinh& $3$ & $12$ & $15$ & $24$ & $2$  \\
		\hline
	\end{longtable}
	Tính tứ phân vị thứ nhất của mẫu số liệu ghép nhóm ở trên (làm tròn kết quả tới hàng phần mười).\\
	\shortans[oly]{$15{,}3$}
	\loigiai{
		Cỡ mẫu là $n=56$.\\
		Tứ phân vị thứ nhất $Q_1$ là $\dfrac{x_{14}+x_{15}}{2}$. \\
		Do $x_{14}$, $x_{15}$ đều thuộc nhóm $[12{,}5 ; 15{,}5)$ nên nhóm này chứa $Q_1$.\\
		Do đó $Q_1=12{,}5+\dfrac{\frac{56}{4}-3}{12} \cdot 3\approx 15{,}3$.
	}
\end{ex}

\begin{ex}
	Cho mẫu số liệu ghép nhóm về điểm số và số học sinh như bảng sau: 
	
	\begin{center}\begin{tabular}{|c|c|c|c|c|c|c|}
			\hline
			Điểm số   & [1 ; 3) & [3 ; 5) & [5 ; 7) & [7 ; 9) & [9 ; 11)\\  
			\hline 
			Số học sinh & 7 & 2 & 4 & 9 & 9 \\ 
			\hline 
		\end{tabular}
	\end{center}
Tìm tứ phân vị thứ ba ${Q_3}$ của mẫu số liệu ghép nhóm đã cho (kết quả làm tròn đến hàng phần trăm).\\
\shortans[oly]{$9{,}28$}
	\loigiai{ 
		
		Bước 1: Tổng tần số là: $N=31$.
		
		Bước 2: Xác định vị trí của $Q_3$: $Q_3$ nằm ở vị trí $\dfrac{3.31}{4}=23.3$.
		
		Bước 3: Xác định lớp chứa $Q_3$: tính tần số tích lũy từ lớp đầu tiên đến khi đạt hoặc vượt qua vị trí của $Q_3$ ta được lớp $[9;11)$.
		
		Bước 4: Xác định các thông số của công thức tính $Q_3$.
		
		Cận dưới của lớp chứa $Q_3$: $L=9$
		
		Tổng tần số của các lớp trước lớp chứa $Q_3$: $F=22$
		
		Tần số của lớp chứa $Q_3$: $f=9$.
		
		Độ rộng lớp chứa $Q_3$: $h=11 - 9=2$.
		
		Áp dụng công thức: $Q_3=L+\left(\dfrac{ \dfrac{3N}{4}-F }{f}\right).h=9+\left(\dfrac{ \dfrac{3.31}{4}-22 }{9}\right).2=9,28$. 
}\end{ex}

\begin{ex}
	Quãng đường (km) các cầu thủ (không tính thủ môn) chạy trong một trận bóng đá tại giải ngoại hạng Anh được cho trong bảng thống kê sau:
\begin{center}
	\begin{tabular}{|c|c|c|c|c|c|}
		\hline Quãng đường & {$[2 ; 4)$} & {$[4 ; 6)$} & {$[6 ; 8)$} & {$[8 ; 10)$} & {$[10 ; 12)$} \\
		\hline Số cầu thủ & 2 & 5 & 6 & 9 & 3 \\
		\hline
	\end{tabular}
\end{center}
	Tìm $a$ sao cho có $25 \%$ số cầu thủ tham gia trận đấu chạy ít nhất $a~(\mathrm{km})$  (kết quả làm tròn đến hàng phần trăm).\\
	\shortans[oly]{$9{,}28$}
\loigiai{
	Số $a$ thỏa mãn có $25 \%$ số cầu thủ tham gia trận đấu chạy ít nhất $a$ (km).\\
	Do đó, $a$ chính là tứ phân vị thứ ba của mẫu số liệu trên.\\
	Cỡ mẫu $n=25$.\\
	Gọi $x_1, x_2 \ldots, x_{25}$ là quãng đường chạy của $25$ cầu thủ và giả sử dãy này đã được sắp xếp theo thứ tự không giảm. Khi đó tứ phân vị thứ ba là $\dfrac{x_{19}+x_{20}}{2}$.\\ Do $x_{19}, x_{20}$ đều thuộc nhóm $[8 ; 10)$ nên nhóm này chứa tứ phân vị thứ ba.\\ Do đó $a=Q_3=8+\dfrac{\dfrac{3.25}{4}-(2+5+6)}{9} \cdot(10-8) \approx 9{,}28$. 
}
\end{ex}
\Closesolutionfile{ans}


